\documentclass[10pt,a4paper]{amsart}
\usepackage[margin=19mm]{geometry}
\usepackage{amsmath,amssymb,mathtools}
\usepackage[scr=rsfs,cal=euler]{mathalfa}
\usepackage{xfrac}
\usepackage[unicode,hidelinks,bookmarks=false]{hyperref}
\newcommand{\ie}{i.e.\ }
\newcommand{\cf}{cf.\ }
\newcommand{\resp}{respectively\ }
\newcommand{\Subsection}{\subsection}
\providecommand{\nobreakdash}{\nobreak\hbox{-}}
\setlength{\emergencystretch}{2em}
\multlinegap0pt
\usepackage[shortlabels]{enumitem}
\usepackage{amssymb}
\usepackage{mathtools}
\usepackage{xcolor}
\usepackage{float}
\usepackage{tikz}
\newtheorem{proposition}{Proposition}
\newcommand{\C}{\mathbb C}
\newcommand{\N}{\mathbb N}
\newcommand{\T}{\mathbb T}
\newcommand{\distT}{d_{\mathbb T}}
\newcommand{\eps}{\varepsilon}
\newcommand{\spancl}[1]{\overline{\operatorname{span}}\{#1\}}
\title{Mathematical Discovery in the Wild:\\ AI-Guided Proofs in Banach Space Theory}
\author{Antonio Acuaviva, Pablo Acuaviva}
\date{Source: arXiv:2607.17388v1 (CC BY 4.0)}
\begin{document}
\maketitle

\noindent\emph{Source and licence.} Mathematical Discovery in the Wild:\\ AI-Guided Proofs in Banach Space Theory. Version arXiv:2607.17388v1 (CC BY 4.0), distributed by arXiv under the Creative Commons Attribution 4.0 International licence, http://creativecommons.org/licenses/by/4.0/, as recorded in the arXiv licence metadata for this exact version and shown on the abstract page https://arxiv.org/abs/2607.17388v1. Full licence notice: \texttt{licenses/arxiv-2607.17388-CC-BY-4.0.txt}.\par\smallskip

\noindent\textit{Editorial scope.} Counted unit: the article's proposition p1-first:prop:rapid (source0.tex lines 773--787). The article's complete original proof of it is delivered with it (source0.tex lines 789--1005, one \texttt{proof} environment of 217 lines). No mathematics is written, repaired or re-derived: the statement and the proof are the article's own lines, in the article's own notation, and no retained line is substituted or re-flowed.  No further source-local context is needed: every cross-reference of every retained line resolves inside the two delivered spans.  Nothing was omitted from any retained span, so the package carries no omission record.  No retained line contains a citation, so the wrapper carries no bibliography and no bibliography entry.  No retained line contains a figure environment, an image-inclusion command or drawing code, so no image is delivered and the package declares no asset dependency.  Editorial formatting only, in addition: an AMS \texttt{amsart} preamble that restates the article's own declarations and macros used by the retained lines, namely the environments \texttt{proposition} on separate counters, together with the standard packages those lines need.  The retained environments print in this document as Proposition 1 in order of appearance, whereas the article prints the same items with the numbering of its own full document.  The complete native source of the article as distributed in the arXiv e-print for this exact version is bundled unchanged as \texttt{sources/205566/source0.tex}.
\begin{proposition}[Rapid flat-block alternative]
\label{p1-first:prop:rapid}
Let $(e_n)_{n=1}^{\infty}$ be a normalized asymptotically monotone basic sequence. Suppose that, for every $\eps>0$, there is no sequence $(x_n)_{n=1}^{\infty}\subset \spancl{e_n:n\in\N}$ of unit vectors satisfying
\begin{equation*}
        \distT(x_n,x_m)>1+\eps \qquad(n,m\in\N,\ n\ne m).
\end{equation*}
Then there are successive flat blocks $(d_j)_{j=1}^{\infty}$ and phases $(\alpha_j)_{j=1}^{\infty}\subset\T$ such that
\begin{equation*}
        \|d_j\|\longrightarrow 1
\end{equation*}
and
\begin{equation*}
        \sup_{N\in\N}\left\|\sum_{j=1}^{N}\alpha_jd_j\right\|<\infty.
\end{equation*}
\end{proposition}

\begin{proof}
Choose two positive null sequences $(\eta_j)_{j=1}^{\infty}$ and $(\gamma_j)_{j=1}^{\infty}$ such that
\begin{equation*}
        0<\eta_j<2^{-j-10}
        \qquad\text{and}\qquad
        0<\gamma_j<2^{-j-10}
        \qquad(j\in\N).
\end{equation*}
For $j\in\N$, set
\begin{equation*}
        \rho_j=\sum_{k=j}^{\infty}(2\eta_k+\gamma_k).
\end{equation*}

By choosing the terms recursively and sufficiently small at each step, we may assume that
\begin{equation*}
        \rho_{j+1}\leq \eta_j \qquad(j\in\N).
\end{equation*}

We first construct, by induction on $j\in\N$, integers $m_j$ and finite families
\begin{equation*}
        B_j=\{b_{j,1},\ldots,b_{j,r_j}\}
\end{equation*}
of $\eta_j$-flat unit blocks, satisfying the following properties:

\begin{enumerate}[label=\textup{(\roman*)}]
\item For every $j\in\N$, the tail after $m_j$ is $(1+\gamma_j)$-monotone: whenever $m_j\leq p\leq q\leq r$ and $a_p,\ldots,a_r\in\C$,
\begin{equation}
\label{p1-first:eq:tail-projection}
        \left\|\sum_{i=p}^{q}a_ie_i\right\|\leq (1+\gamma_j)\left\|\sum_{i=p}^{r}a_ie_i\right\|.
\end{equation}

\item\label{p1-first:item:maximal-family} For every $j\in\N$, the family $B_j$ is a finite maximal family of $\eta_j$-flat unit blocks supported after $m_j$ such that
\begin{equation*}
        \distT(b_{j,i},b_{j,\ell})>1+\eta_j \qquad(1\leq i<\ell\leq r_j).
\end{equation*}

\item The families are successive: every member of $B_j$ is supported before $m_{j+1}$.
\end{enumerate}

Here, a vector is supported after $m_j$ if its support is contained in $\{m_j+1,m_j+2,\ldots\}$. Maximality in \ref{p1-first:item:maximal-family} is meant with respect to inclusion among families satisfying the displayed separation condition. Thus, once $B_j$ has been chosen, no further $\eta_j$-flat unit block supported after $m_j$ can be added to $B_j$ while preserving pairwise toroidal separation greater than $1+\eta_j$. Equivalently, every $\eta_j$-flat unit block supported after $m_j$ lies within distance at most $1+\eta_j$, up to multiplication by a unimodular scalar, of some member of $B_j$.

We now carry out the construction. By asymptotic monotonicity, for each $j\in\N$ there is an integer $N_j$ such that \eqref{p1-first:eq:tail-projection} holds whenever $N_j\leq p\leq q\leq r$.

Choose $m_1\geq N_1$. Having chosen $m_j$, consider the class of all $\eta_j$-flat unit blocks supported after $m_j$. By the standing assumption, this class contains no infinite subfamily whose distinct members are pairwise at $\distT$-distance greater than $1+\eta_j$. Hence a greedy selection process gives a finite maximal family
\begin{equation*}
        B_j=\{b_{j,1},\ldots,b_{j,r_j}\}
\end{equation*}
with the separation property in \ref{p1-first:item:maximal-family}. For each $1\leq i\leq r_j$, choose a flat block $d_{j,i}$ such that
\begin{equation}
\label{p1-first:eq:representing-block}
        b_{j,i}=\frac{d_{j,i}}{\|d_{j,i}\|}, \qquad |\|d_{j,i}\|-1|\leq \eta_j.
\end{equation}
Since $B_j$ is finite, we may choose $m_{j+1}\geq N_{j+1}$ larger than every index appearing in the supports of the vectors in $B_j$. This defines $B_j$ and $m_{j+1}$, and hence completes the recursive construction.

The maximality of $B_j$ gives the following covering property. If $u$ is any $\eta_j$-flat unit block supported after $m_j$, then there are $1\leq i\leq r_j$ and $\theta\in\T$ such that
\begin{equation}
\label{p1-first:eq:covering}
        \|b_{j,i}-\theta u\|\leq 1+\eta_j.
\end{equation}
Indeed, otherwise $u$ could be added to $B_j$, contradicting maximality.

This finishes the construction of the finite covers $B_j$. We now use them to build finite flat block combinations with uniformly bounded partial sums. The construction is done first at a fixed finite length $L$, and then a diagonal argument will let $L\to\infty$.

Fix $L\in\N$. We shall construct, by backward induction on $j=L,L-1,\ldots,1$, indices
\begin{equation*}
        \iota_j^{(L)}\in\{1,\ldots,r_j\},
\end{equation*}
phases $\theta_j^{(L)}\in\T$ for $1\leq j<L$, and flat blocks
\begin{equation*}
        z_L^{(L)},z_{L-1}^{(L)},\ldots,z_1^{(L)}
\end{equation*}
such that the following hold.

\begin{enumerate}[label=\textup{(\alph*)}]
\item\label{p1-first:item:last-level} At the last level,
\begin{equation*}
        z_L^{(L)}=d_{L,\iota_L^{(L)}}.
\end{equation*}

\item\label{p1-first:item:backward-recursion} For every $1\leq j<L$, if
\begin{equation*}
        u_{j+1}^{(L)}=\frac{z_{j+1}^{(L)}}{\|z_{j+1}^{(L)}\|},
\end{equation*}
then
\begin{equation}
\label{p1-first:eq:covering-step}
        \left\|b_{j,\iota_j^{(L)}}-\theta_j^{(L)}u_{j+1}^{(L)}\right\|\leq 1+\eta_j,
\end{equation}
and
\begin{equation}
\label{p1-first:eq:z-recursion}
        z_j^{(L)}=d_{j,\iota_j^{(L)}}-\theta_j^{(L)}z_{j+1}^{(L)}.
\end{equation}

\item\label{p1-first:item:norm-control} For every $1\leq j\leq L$,
\begin{equation}
\label{p1-first:eq:z-error-target}
        |\|z_j^{(L)}\|-1|\leq \rho_j.
\end{equation}
\end{enumerate}

We start at level $L$. Choose any $\iota_L^{(L)}\in\{1,\ldots,r_L\}$ and set
\begin{equation*}
        z_L^{(L)}=d_{L,\iota_L^{(L)}}.
\end{equation*}
This gives \ref{p1-first:item:last-level}. Moreover, by \eqref{p1-first:eq:representing-block},
\begin{equation*}
        |\|z_L^{(L)}\|-1|\leq \eta_L\leq \rho_L,
\end{equation*}
so \ref{p1-first:item:norm-control} holds at level $L$.

Now suppose that $1\leq j<L$ and that the construction has been completed at level $j+1$. In particular, \ref{p1-first:item:norm-control} gives
\begin{equation*}
        |\|z_{j+1}^{(L)}\|-1|\leq \rho_{j+1}\leq \eta_j.
\end{equation*}
Since the families $B_1,B_2,\ldots$ are successive, $z_{j+1}^{(L)}$ is supported after $m_j$. Hence the normalized vector
\begin{equation*}
        u_{j+1}^{(L)}=\frac{z_{j+1}^{(L)}}{\|z_{j+1}^{(L)}\|}
\end{equation*}
is an $\eta_j$-flat unit block supported after $m_j$. By the covering property \eqref{p1-first:eq:covering}, there are $\iota_j^{(L)}\in\{1,\ldots,r_j\}$ and $\theta_j^{(L)}\in\T$ such that \eqref{p1-first:eq:covering-step} holds. We then define
\begin{equation*}
        z_j^{(L)}=d_{j,\iota_j^{(L)}}-\theta_j^{(L)}z_{j+1}^{(L)}.
\end{equation*}
Thus \ref{p1-first:item:backward-recursion} holds at level $j$. Since the two terms have successive disjoint supports and unimodular coefficients on their supports, $z_j^{(L)}$ is again a flat block.

It remains to verify \ref{p1-first:item:norm-control} at level $j$. Put
\begin{equation*}
        s_j^{(L)}=\|d_{j,\iota_j^{(L)}}\|,
        \qquad
        t_{j+1}^{(L)}=\|z_{j+1}^{(L)}\|.
\end{equation*}
By \eqref{p1-first:eq:representing-block} and the induction hypothesis,
\begin{equation*}
        |s_j^{(L)}-1|\leq\eta_j,
        \qquad
        |t_{j+1}^{(L)}-1|\leq\rho_{j+1}.
\end{equation*}
By \ref{p1-first:item:backward-recursion}, and by the definitions of $s_j^{(L)}$, $t_{j+1}^{(L)}$ and $u_{j+1}^{(L)}$, we have
\begin{equation*}
        z_j^{(L)}=s_j^{(L)}b_{j,\iota_j^{(L)}}-\theta_j^{(L)}t_{j+1}^{(L)}u_{j+1}^{(L)}.
\end{equation*}
Therefore,
\begin{equation}
\label{p1-first:eq:upper-z}
\begin{aligned}
        \|z_j^{(L)}\|
        &=\left\|s_j^{(L)}b_{j,\iota_j^{(L)}}-\theta_j^{(L)}t_{j+1}^{(L)}u_{j+1}^{(L)}\right\|\\
        &\leq \left\|b_{j,\iota_j^{(L)}}-\theta_j^{(L)}u_{j+1}^{(L)}\right\|+|s_j^{(L)}-1|+|t_{j+1}^{(L)}-1|\\
        &\leq 1+\eta_j+\eta_j+\rho_{j+1}\\
        &=1+2\eta_j+\rho_{j+1}.
\end{aligned}
\end{equation}
For the lower estimate, $d_{j,\iota_j^{(L)}}$ is an initial interval projection of $z_j^{(L)}$. Since $m_j\geq N_j$, \eqref{p1-first:eq:tail-projection} gives
\begin{equation*}
        \|d_{j,\iota_j^{(L)}}\|\leq (1+\gamma_j)\|z_j^{(L)}\|.
\end{equation*}
Therefore
\begin{equation}
\label{p1-first:eq:lower-z}
        \|z_j^{(L)}\|\geq \frac{1-\eta_j}{1+\gamma_j}\geq 1-\eta_j-\gamma_j,
\end{equation}
where in the last inequality we used the elementary estimate
\begin{equation*}
        \frac{1-a}{1+b}\geq 1-a-b \qquad(a,b>0).
\end{equation*}
Combining \eqref{p1-first:eq:upper-z} and \eqref{p1-first:eq:lower-z}, we obtain
\begin{equation*}
        |\|z_j^{(L)}\|-1|\leq 2\eta_j+\gamma_j+\rho_{j+1}=\rho_j.
\end{equation*}
Thus \ref{p1-first:item:norm-control} holds at level $j$. This proves the induction step and hence completes the backward construction.

Unwinding \eqref{p1-first:eq:z-recursion}, and setting
\begin{equation*}
        \alpha_1^{(L)}=1, \qquad \alpha_j^{(L)}=(-1)^{j-1}\theta_1^{(L)}\cdots\theta_{j-1}^{(L)} \quad(2\leq j\leq L),
\end{equation*}
we obtain
\begin{equation}
\label{p1-first:eq:finite-expansion}
        z_1^{(L)}=\sum_{j=1}^{L}\alpha_j^{(L)}d_{j,\iota_j^{(L)}}.
\end{equation}
In particular, by \eqref{p1-first:eq:z-error-target},
\begin{equation}
\label{p1-first:eq:z1-bound}
        \|z_1^{(L)}\|\leq 1+\rho_1 \qquad(L\in\N).
\end{equation}

Thus, for each fixed $L\in\N$, the construction gives a finite sequence
\begin{equation*}
        d_{1,\iota_1^{(L)}}<d_{2,\iota_2^{(L)}}<\ldots<d_{L,\iota_L^{(L)}}
\end{equation*}
and phases $\alpha_1^{(L)},\ldots,\alpha_L^{(L)}\in\T$. If $N\leq L$, then the partial sum over the first $N$ levels is an initial interval projection of $z_1^{(L)}$. Applying \eqref{p1-first:eq:tail-projection} at level $1$, we get
\begin{equation}
\label{p1-first:eq:finite-partial-bound}
        \left\|\sum_{j=1}^{N}\alpha_j^{(L)}d_{j,\iota_j^{(L)}}\right\|\leq (1+\gamma_1)(1+\rho_1)=:C.
\end{equation}


Finally, we pass from finite lengths to an infinite sequence. For each fixed $j\in\N$ and $L \geq j$, the indices $\iota_j^{(L)}$ take values in the finite set $\{1,\ldots,r_j\}$, while the phases $\alpha_j^{(L)}$ lie in the compact set $\T$. A diagonal compactness argument gives a subsequence $L_k\to\infty$ such that, for each fixed $j\in\N$,
\begin{equation*}
        \iota_j^{(L_k)}=\iota_j \text{ eventually},
        \qquad
        \alpha_j^{(L_k)}\to\alpha_j\in\T.
\end{equation*}
Set
\begin{equation*}
        d_j=d_{j,\iota_j}\qquad(j\in\N).
\end{equation*}
Then $d_1<d_2<\ldots$, and \eqref{p1-first:eq:representing-block} gives $\|d_j\|\to1$. Fixing $N\in\N$ in \eqref{p1-first:eq:finite-partial-bound} and passing to the limit along $L_k$ yields
\begin{equation*}
        \left\|\sum_{j=1}^{N}\alpha_jd_j\right\|\leq C.
\end{equation*}
Since $N\in\N$ was arbitrary,
\begin{equation*}
        \sup_{N\in\N}\left\|\sum_{j=1}^{N}\alpha_jd_j\right\|<\infty.
\end{equation*}
This proves the proposition.
\end{proof}

\end{document}
